% (a) 贯穿图 V={1,2,3,4}, E={12,13,23,24,34}：顶点2的邻域 N(2)={1,3,4}（绿色），连接2的三条边加粗；各顶点旁标度 2,3,3,2，度数和10=2*5。
% (b) 有向图 1->2, 3->2, 2->4：N^-(2)={1,3}（黄色），N^+(2)={4}（绿色），入度2、出度1。
\begin{tikzpicture}[>=Stealth,
  v/.style={circle,draw=FigureStroke,fill=white,line width=.8pt,minimum size=6.5mm,inner sep=0pt,text=FigureInk},
  cv/.style={v,fill=FigureBlueFill},
  nv/.style={v,fill=FigureGreenFill},
  iv/.style={v,fill=FigureYellowFill},
  e/.style={draw=FigureStroke,line width=.8pt},
  he/.style={draw=FigureInk,line width=1.8pt},
  a/.style={draw=FigureStroke,line width=.8pt,->,shorten >=1pt},
  ha/.style={draw=FigureInk,line width=1.5pt,->,shorten >=1pt},
  deg/.style={font=\FigureNoteFont\fontsize{7.5}{9}\selectfont,text=FigureMuted},
  note/.style={font=\FigureNoteFont\fontsize{7.5}{9}\selectfont,text=FigureMuted,align=center}]
\begin{scope}
  \node[nv] (1) at (-1.2,0) {$1$}; \node[cv] (2) at (0,.85) {$2$};
  \node[nv] (3) at (0,-.85) {$3$}; \node[nv] (4) at (1.2,0) {$4$};
  \draw[e] (1)--(3); \draw[e] (3)--(4);
  \draw[he] (2)--(1); \draw[he] (2)--(3); \draw[he] (2)--(4);
  \node[deg,left] at (1.west) {$d_1=2$};
  \node[deg,above] at (2.north) {$d_2=3$};
  \node[deg,below] at (3.south) {$d_3=3$};
  \node[deg,right] at (4.east) {$d_4=2$};
  \node at (0,-2.0) {(a) $N(2)=\{1,3,4\}$，$d_2=3$};
  \node[note] at (0,-2.55) {度数和$2+3+3+2=10=2m$};
\end{scope}
\begin{scope}[xshift=4.6cm]
  \node[iv] (1) at (-1.2,.6) {$1$}; \node[iv] (3) at (-1.2,-.6) {$3$};
  \node[cv] (2) at (0,0) {$2$}; \node[nv] (4) at (1.2,0) {$4$};
  \draw[ha] (1)--(2); \draw[ha] (3)--(2); \draw[ha] (2)--(4);
  \node at (0,-2.0) {(b) $N^-(2)=\{1,3\}$，$N^+(2)=\{4\}$};
  \node[note] at (0,-2.55) {入度$d_2^-=2$，出度$d_2^+=1$};
\end{scope}
\end{tikzpicture}
